\documentclass[book.tex]{subfiles}

\begin{document}
\chapter{Symmetry Groups}

\label{chap:lorentz_group}
\noindent\rule{11cm}{0.4pt} \\
\textbf{Prerequisites:} \autoref{chap:lie} \\
\textbf{Difficulty Level:} ** \\
\noindent\rule{11cm}{0.4pt}
\vspace{5mm}

In this chapter, we introduce a number of symmetry groups that arise in physical systems. All the groups we introduce here are examples of topological groups, where they inherit a topological structure from a larger space (typically $\mathbb{R}^n$ or $\mathbb{R}^{n \times n}$). In fact, as we will learn in a later chapter, these groups also inherit a differentiable structure from the larger space, making them all examples of "Lie Groups" (or differentiable groups). 

\section{The Orthogonal Groups $O(N)$}

The orthogonal group $O(3)$ is the collection of matrices with orthogonal columns in $R^3$. More general, $O(N)$ is the collection of $N \times N$ matrices which have orthogonal columns.

\begin{align}
O(N) &= \{ Q \in \mathbb{R}^{N \times N} | Q^T Q = QQ^T = I\}
\end{align}


\section{The Rotation Groups $SO(N)$}

The rotation group $SO(3)$ is the collection of three dimensional rotations, and is denoted $\textrm{SO}(3)$. $SO(3)$ is related to $O(3)$ since it is the subset of $O(3)$ with determinant $1$. Mathematically, we say that $\det$ is a homomorphism from $O(3)$ to the real numbers minus $0$.
\begin{align}
\det : O(N) \to \mathbb{R} \setminus \{0 \}
\end{align}
This fact holds since we have that
\begin{align}
\det(AB) &= \det A \det B
\end{align}
$SO(3)$ is the subset of $O(3)$ that maps to $1$ under this transformation. An alternative way of describing $SO(3)$ is that it is the collection of orthogonal matrices whose columns are orthonormal (and not just orthogonal).

All elements of $SO(3)$ are length preserving since rotations don't change the length of vectors. This holds true more generally for $SO(N)$ and can be shown using the fact that column vectors are orthonormal. Another important fact is that rotations are linear. That is, if $R$ is a rotation and $u$ and $v$ are vectors, we have
\begin{align}
R \cdot (u + v) &= R \cdot u + R \cdot v
\end{align}
For this reason, $SO(3)$ can be identified with the collection of real valued $3x3$ matrices with unit determinant.


\section{The Euclidean Group $E(N)$}

The Euclidean group $E(n)$ denotes all the distance preserving transformations (isometries) on $\mathbb{R}^n$. This collection includes all translations, rotations, and reflections. The group $E(3)$ commonly arises when considering point clouds or collections of atoms in a molecule or material.

\section{The Galilean Group}

The Galilean group consists of the collection of transformations between two reference frames which only differ by constant relative motion with respect to one another. Consider $(x, t)$ where $x \in \mathbb{R}^3$ is a coordinate and $t \in \mathbb{R}$ is time. Then uniform motion is represented by the transformation
\begin{align}
(x, t) \mapsto (x + tv, t)
\end{align}
Rotation is represented by the transformation
\begin{align}
(x, t) \mapsto (R x, t)
\end{align}
where $R$ is a rotation matrix. Translations are represented by
\begin{align}
(x, t) \mapsto (x + v, t + s)
\end{align}
The group of Galilean transformations consists of all the compositions of these three types of transformations. 

\section{The Lorentz Group $O(1, 3)$}

The Lorentz group is the group of all Lorentz transformations defined on Minkowski space. Many fundamental equations in physics are invariant under Lorentz transformations. Mathematically, we write the Lorentz group as $O(1, 3)$ the set of orthogonal transformations on $\mathbb{R}^4$ that preserve the quadratic form
\begin{align}
    (t, x, y, z) \mapsto t^2 - x^2 - y^2 - z^2 
\end{align}
The Lorentz group is often studied through its Lie group $\mathfrak{so}(1, 3)$
\begin{align}
    \mathfrak{so}(1, 3) &= \{ X \in \mathbb{R}^{4 \times 4} | e^{tX} \in SO(1, 3), t \in \mathbb{R} \}
\end{align}

\section{Exercises}
\begin{enumerate}
    \item Prove that the Galilean group has dimension 10.
    \item Prove that $SO(N)$ preserves the norm of vectors in $\mathbb{R}^N$
\end{enumerate}

\end{document}